% !TEX root =../ECE343_Slides_Rogers.tex

%%%%%%%%%%%%%%%%%%%% LO 3.3%%%%%%%%%%%%%%%%%
%%%%%%%%%%%%%%%%%%%%%%%%%%%%%%%%%%%%%%%%%%%%%%%%%%%%%%
%%%%%%%%%%%%%%%%%%%%%%%%%%%%%%%%%%%%%%%%%%%%%%%%%%%%%%

\section{Schedule}
\renewcommand{\arraystretch}{1}
\begin{frame}{}
	\begin{textblock}{14}(0.5,1.75)
	\def\myfontsize{footnotesize}
	
\setcounter{mylnum}{1}
\begin{center}
\begin{\myfontsize}
\begin{tabular}{| c | l | c | l |}
\hline
\textbf{Lesson} 		  	& \textbf{Topics}             & \pbox{5cm}{\textbf{Reading}}	& \textbf{Notes}		\\ \hline\hline
\ifnum\lessonNum=\value{mylnum} \rowcolor{green!20!white}\fi
	\themylnum  
	& Course Overview  \& Introduction		
	& Ch.1 
	&  LO 1.1-1.2 
	\stepcounter{mylnum}\\ \hline \hline
\rowcolor{gray!20!white}\multicolumn{4}{| c |}{\textbf{Chapter 2: Electric and Magnetic Fields}}	\\ \hline\hline
\ifnum\lessonNum=\value{mylnum} \rowcolor{green!20!white}\fi
	\themylnum 	
	& Fields, Parameters, Field Quantities	
	& Ch. 2	
	& LO 2.1-2.2			
\stepcounter{mylnum}\\ \hline\hline
\rowcolor{gray!20!white}\multicolumn{4}{| c |}{\textbf{Chapter 3: Transmission Lines}}	 \\ \hline\hline	
\ifnum\lessonNum=\value{mylnum} \rowcolor{green!20!white}\fi
	\themylnum 	
	& Lumped Element Model				
	& 3.1-3.4 
	& LO 3.1		
\stepcounter{mylnum}\\ \hline
\ifnum\lessonNum=\value{mylnum} \rowcolor{green!20!white}\fi
	\themylnum 	
	& Telegrapher's Equations
	& 3.5-3.6
	&  LO 3.1	
\stepcounter{mylnum}\\ \hline
\ifnum\lessonNum=\value{mylnum} \rowcolor{green!20!white}\fi
	\themylnum 	
	&  Characteristic Impedance \& Wave Propagation	
	& 3.7-3.9		
	& LO 3.2						
\stepcounter{mylnum}\\ \hline
\ifnum\lessonNum=\value{mylnum} \rowcolor{green!20!white}\fi
	\themylnum 	
	& Coaxial,  Reflection Coefficient, VSWR	
	& 3.10,3.12-3.14	
	& LO 3.2-3.3		
\stepcounter{mylnum}\\ \hline 
\ifnum\lessonNum=\value{mylnum} \rowcolor{green!20!white}\fi
	\themylnum 	
	&  Terminated Transmission Lines	
	&  3.15,3.17-3.19
	& LO 3.4		
\stepcounter{mylnum}\\ \hline 
\ifnum\lessonNum=\value{mylnum} \rowcolor{green!20!white}\fi
	\themylnum	
	& Impedance Matching	
	& 3.21-3.23
	& LO 3.4 			
\stepcounter{mylnum}\\ \hline 
\ifnum\lessonNum=\value{mylnum} \rowcolor{green!20!white}\fi
	\themylnum  
	& Impedance Matching	
	& 3.21-3.23
	& LO 3.4		
\stepcounter{mylnum}\\ \hline 
\ifnum\lessonNum=\value{mylnum} \rowcolor{green!20!white}\fi
	\themylnum	
	& Smith Charts
	& Supplemental
	& LO 3.5
\stepcounter{mylnum}\\ \hline 
\ifnum\lessonNum=\value{mylnum} \rowcolor{green!20!white}\fi
	\themylnum	
	& Smith Charts
	& Supplemental 
	& LO 3.5					
\stepcounter{mylnum}\\ \hline 
\ifnum\lessonNum=\value{mylnum} \rowcolor{green!20!white}\fi
	\themylnum	
	& Smith Charts
	& Supplemental 
	& LO 3.5	
\stepcounter{mylnum}\\ \hline 
\ifnum\lessonNum=\value{mylnum} \rowcolor{green!20!white}\fi
	\themylnum	
	& Smith Charts
	& Supplemental 
	& LO 3.5	
\stepcounter{mylnum}\\ \hline 

\ifnum\lessonNum=\value{mylnum} \rowcolor{green!20!white}\fi
	\themylnum	
	& Transients		
	& Supplemental 
	& LO 3.6	 
\stepcounter{mylnum}\\ \hline \hline
\ifnum\lessonNum=\value{mylnum} \rowcolor{green!20!white}\fi
	\themylnum 	
	& Flex Day
	& 
	& 	
\stepcounter{mylnum}\\ \hline 
\ifnum\lessonNum=\value{mylnum} \rowcolor{green!20!white}\fi
	\themylnum 	
	& Reflection/Transmission Coefficient Lab
	& Lab Packet
	& LO 3.7		
\stepcounter{mylnum}\\ \hline 
\ifnum\lessonNum=\value{mylnum} \rowcolor{green!20!white}\fi
	\themylnum 	
	& Reflection/Transmission Coefficient Lab
	& Lab Packet
	& LO 3.7		
\stepcounter{mylnum}\\ \hline 
\rowcolor{gray!20!white}\multicolumn{4}{| c |}{\textbf{Chapter 4: Vector Analysis}} \\ \hline\hline
\ifnum\lessonNum=\value{mylnum} \rowcolor{green!20!white}\fi
	\themylnum	
	& Vector \& Coordinate System Review
	& Ch. 4
	&LO 4.1-4.4 \textbf{GR 1 due}	
\stepcounter{mylnum}\\ \hline 
\rowcolor{gray!20!white}\multicolumn{4}{| c |}{\textbf{Chapter 5: Electrostatics}}\\ \hline\hline
\ifnum\lessonNum=\value{mylnum} \rowcolor{green!20!white}\fi
	\themylnum	
	& Coulomb's Law
	& 5.1-5.3
	& LO 5.1
\stepcounter{mylnum}\\ \hline 
\ifnum\lessonNum=\value{mylnum} \rowcolor{green!20!white}\fi
	\themylnum	
	& Surface and Volume Charge			
	& 5.4
	& LO 5.1
\stepcounter{mylnum}\\ \hline 
\ifnum\lessonNum=\value{mylnum} \rowcolor{green!20!white}\fi
	\themylnum	
	& Gauss' Law (Integral Form)
	& 5.5-5.6
	& LO 5.2		 					
\stepcounter{mylnum}\\ \hline 
\ifnum\lessonNum=\value{mylnum} \rowcolor{green!20!white}\fi
	\themylnum	
	& Gauss' Law (Differential Form)		
	& 5.7-5.8
	& LO 5.2	 	
\stepcounter{mylnum}\\ \hline 
\ifnum\lessonNum=\value{mylnum} \rowcolor{green!20!white}\fi
	\themylnum	
	& Electric Potential	
	& 5.9-5.10, 5.15
	&  LO 5.3	
\stepcounter{mylnum}\\ \hline 
% <---- line 129 should be 20 lessons, based on 6 lines per lesson, starting at line 9 
\end{tabular}			
\end{\myfontsize}
\end{center}

	\end{textblock}
\end{frame}

%%%%%%%%%%%%%%%%%%%%%%%%%%%%%%%%%%%%%%%%%%%%%%%%%%%%%
%%%%%%%%%%%%%%%%%%%%%%%%%%%%%%%%%%%%%%%%%%%%%%%%%%%%%

\def\obja{I can describe the origin of the transmission line equation, understand the method of developing a solution, and interpret the solution.}
\def\objb{I can explain the physical meaning of and calculate the following quantities: propagation constant, phase constant, attenuation constant, characteristic impedance, VSWR, reflection coefficient, and transmission coefficient.}
\def\objc{I can use the transmission line equation to calculate the voltage/current at an arbitrary point on a specified transmission line.}
\def\objd{ I can use the transmission line equation to match a load to a source using a length of transmission line.}
\def\obje{ I can identify the various quantities represented on the Smith Chart.}
\def\objf{ I can use a Smith Chart to design a matching network for a given source and load, using series components, parallel components, shorted stubs, and open stubs.}
\def\objg{ I can plot the transient response of a load connected to a source through a transmission line. }
\def\objh{ I can measure the reflection coefficient across a range of frequencies and interpret the measurement using a network analyzer.}
\begin{frame}{Ch 3 Learning Objectives}
	\begin{textblock}{14}(0.5,3)
	\begin{scriptsize}
	\begin{enumerate}[{3}.1]
		\ifnum\objNum=1
			\item \alert{ \obja }
		\else 
			\item \obja
		\fi 
		\ifnum\objNum=23
			\item \alert{ \objb }
			\item \alert{ \objc }
		\else 
			\item \objb
			\item \objc
		\fi 
%		\ifnum\objNum=3
%			\item \alert{ \objc }
%		\else 
%			\item \objc
%		\fi 
		\ifnum\objNum=4
			\item \alert{ \objd }
		\else 
			\item \objd
		\fi 
		\ifnum\objNum=56
			\item \alert{ \obje }
		\else 
			\item \obje
		\fi 
		\ifnum\objNum=56
			\item \alert{ \objf }
		\else 
			\item \objf
		\fi 
		\ifnum\objNum=7
			\item \alert{ \objg }
		\else 
			\item \objg
		\fi 
		\ifnum\objNum=8
			\item \alert{ \objh }
		\else 
			\item \objh
		\fi 

	\end{enumerate}
	\end{scriptsize}
	\end{textblock}
\end{frame}

%%%%%%%%%%%%%%%%%%%%%%%%%%%%%%%%%%%%%%%%%%%%%%%%%%%%%
%%%%%%%%%%%%%%%%%%%%%%%%%%%%%%%%%%%%%%%%%%%%%%%%%%%%%

\section{Transients}

\begin{frame}{Transients}
	\begin{textblock}{8}(0,2)
		\begin{myblock}[0.9]{Method}
		\begin{scriptsize}
		\begin{itemize}
			\item At $t=0$, source voltage is divided between \handoutgap{0.5in}{\myz{S}} and \handoutgap{0.5in}{\myz{0}}
			
			\item Wave propagates from \handoutgap{0.5in}{$z=-\ell$} to  \handoutgap{0.5in}{$z=0$} at velocity  \handoutgap{1in}{$u_p$}
			
			\item  Wave \handoutgap{0.5in}{reflects} at boundary of \handoutgap{0.5in}{tx line \& load}
			
			\item Wave propagates from \handoutgap{0.5in}{$z=0$} to  \handoutgap{1in}{$z=-\ell$} at velocity  \handoutgap{0.5in}{$u_p$} 
			
			\item  Wave \handoutgap{0.5in}{reflects} at boundary of \handoutgap{0.5in}{source \& tx line}
			
			\item Voltage converges to \handoutgap{0.5in}{steady state} of \handoutgap{0.5in}{$V_0^+ \lp\dfrac{\myz{L}}{\myz{L}+\myz{s}} \rp$}
		\end{itemize}
		\end{scriptsize}
	\end{myblock}
	\end{textblock}
		\begin{textblock}{8}(8,4)
			\scalebox{0.4}{
				\begin{circuitikz}[scale=1]
					\draw (1.5,0) to[open, *-*] ++(0,2)
										to[short, -*] ++(8.5,0);
					\draw (10,0) to[short,*-] ++(-10,0);
					\draw (0,0) to[short] ++(-2,0)
										to[V,a^=\myv{s}, invert] ++(0,2)
										to[twoport, t=\myz{s}] ++(2,0)
										to[closing switch] ++(2,0); % source
					\draw (10,2) to[short] ++(2,0)
										to[twoport, t=\myz{L}] ++(0,-2)
										to[short] ++(-2,0); % load
					\draw[->] (0.5,3) -++ (9.5,0);
					\node at (5,3.25) {$z$};
					\draw[fill=gray!50!white, thick] (2,1.75) rectangle ++(7,0.5); % top tx line
					\draw[fill=gray!50!white, thick] (2,-0.25) rectangle ++(7,0.5); % bottom tx line
					\draw[->] (1.5,2.25) --++(0,-9) ; % left boundary
					\node at (1.5,2.5){$-\ell$};
					\draw[->] (10,2.25) --++(0,-9); % right boundary
					\node at (10,2.5) {0};
					
					\foreach \y in {1,2,3}{
						\draw[dashed] (-0.5, -2*\y) --++ (10.75,0)  node[right] {\y$t_{\ell}$};
					}
					% red lines for bounces
					\draw[thick, red,->] (1.5,0) --++ (8.5,-2) node[above, midway] {$\myv{0}^+$}; 
					\draw[thick, red,->] (10,-2) --++ (-8.5,-2) node[above, midway] {$\myv{0}^- = \myv{0}^+\Gamma_L$};
					\draw[thick, red,->] (1.5,-4) --++ (8.5,-2)  node[above, midway] {$\myv{0}^{-} \Gamma_s = \myv{0}^+\Gamma_L \Gamma_s$};
				\end{circuitikz}
				} % end scalebox

	\end{textblock}
\end{frame}

%%%%%%%%%%%%%%%%%%%%%%%%%%%%%%%%%%%%%%%%%%%%
%%%%%%%%%%%%%%%%%%%%%%%%%%%%%%%%%%%%%%%%%%%%

\begin{frame}{Transients}
	\begin{textblock}{14.5}(0.5,2.5)
	\begin{myblock}[0.98]{Method}
		\begin{itemize}
			\item At $t=0$, source voltage is divided between \handoutgap{1in}{\myz{S}} and \handoutgap{1in}{\myz{0}}
			
			\item Wave propagates from \handoutgap{1in}{$z=-\ell$} to  \handoutgap{1in}{$z=0$} at velocity  \handoutgap{1in}{$u_p$}
			
			\item  Wave \handoutgap{1in}{reflects} at boundary of \handoutgap{1in}{tx line \& load}
			
			\item Wave propagates from \handoutgap{1in}{$z=0$} to  \handoutgap{1in}{$z=-\ell$} at velocity  \handoutgap{1in}{$u_p$} 
			
			\item  Wave \handoutgap{1in}{reflects} at boundary of \handoutgap{1in}{source \& tx line}
			
			\item Voltage converges to \handoutgap{1in}{steady state} of \handoutgap{1in}{$V_0^+ \lp\dfrac{\myz{L}}{\myz{L}+\myz{s}} \rp$}
			
		\end{itemize}
	\end{myblock}
	\end{textblock}
\end{frame}


%%%%%%%%%%%%%%%%%%%%%%%%%%%%%%%%%%%%%%%%%%%%%%%%%%%%%
%%%%%%%%%%%%%%%%%%%%%%%%%%%%%%%%%%%%%%%%%%%%%%%%%%%%%

\begin{frame}{Transient Example}
			\begin{textblock}{14.5}(0.5,2)
			\begin{center}
			\begin{circuitikz}[scale=0.75]
					\draw (1.5,0) to[open, *-*] ++(0,2)
										to[short, -*] ++(8.5,0);
					\draw (10,0) to[short,*-] ++(-10,0);
					\draw (0,0) to[short] ++(-2,0)
										to[V,a^=10\volts, invert] ++(0,2)
										to[R=25\ohms] ++(2,0)
										to[closing switch, a^=\scriptsize{$t\equals 0$}] ++(2,0); % source
					\draw (10,2) to[short] ++(1,0)
										to[R=150\ohms] ++(0,-2)
										to[short] ++(-1,0); % load
					\draw (1.5,3) --++ (8.5,0) node[above, midway] (5,3.25) {$\ell=400$m};
					\draw(1.5,2.75) --++ (0,0.5);
					\draw(10,2.75) --++ (0,0.5);
					\draw[fill=gray!50!white, thick] (2,1.75) rectangle ++(7,0.5); % top tx line
					\draw[fill=gray!50!white, thick] (2,-0.25) rectangle ++(7,0.5); % bottom tx line
					\draw[->] (1.5,2.25) --++(0,-9) ; % left boundary
					\draw[->] (10,2.25) --++(0,-9); % right boundary
					\node[align=left] at (5,1) {\myz{0}=50\ohms \\ $u_p = 0.67c$};
					
					\foreach \y in {1,2,3}{
						\draw[dashed] (-0.5, -2*\y) --++ (10.75,0)  node[right] {\y$t_{\ell}$};
					}
					% red lines for bounces
					\draw[thick, red,->] (1.5,0) --++ (8.5,-2);% node[above] {$\myv{0}^+$}; 
					\draw[thick, red,->] (10,-2) --++ (-8.5,-2);%node[above] {$\myv{0}^- = \myv{0}^+\Gamma_L$};
					\draw[thick, red,->] (1.5,-4) --++ (8.5,-2);% node[above] {$\myv{0}^{-} \Gamma_s = \myv{0}^+\Gamma_L \Gamma_s$};
				\end{circuitikz}
				\end{center}
			\end{textblock}
\end{frame}

%%%%%%%%%%%%%%%%%%%%%%%%%%%%%%%%%%%%%%%%%%%%%%%%%%%%%
%%%%%%%%%%%%%%%%%%%%%%%%%%%%%%%%%%%%%%%%%%%%%%%%%%%%%

\begin{frame}{Transient Example}
			\begin{textblock}{8}(0.5,2)
			\begin{center}
			\scalebox{0.5}{
			\begin{circuitikz}[scale=0.75]
					\draw (1.5,0) to[open, *-*] ++(0,2)
										to[short, -*] ++(8.5,0);
					\draw (10,0) to[short,*-] ++(-10,0);
					\draw (0,0) to[short] ++(-2,0)
										to[V,a^=10\volts, invert] ++(0,2)
										to[R=25\ohms] ++(2,0)
										to[closing switch, a^=\scriptsize{$t\equals 0$}] ++(2,0); % source
					\draw (10,2) to[short] ++(1,0)
										to[R=150\ohms] ++(0,-2)
										to[short] ++(-1,0); % load
					\draw (1.5,3) --++ (8.5,0) node[above, midway] (5,3.25) {$\ell=400$m};
					\draw(1.5,2.75) --++ (0,0.5);
					\draw(10,2.75) --++ (0,0.5);
					\draw[fill=gray!50!white, thick] (2,1.75) rectangle ++(7,0.5); % top tx line
					\draw[fill=gray!50!white, thick] (2,-0.25) rectangle ++(7,0.5); % bottom tx line
					\draw[->] (1.5,2.25) --++(0,-9) ; % left boundary
					\draw[->] (10,2.25) --++(0,-9); % right boundary
					\node[align=left] at (5,1) {\myz{0}=50\ohms \\ $u_p = 0.67c$};
					
					\foreach \y in {1,2,3}{
						\draw[dashed] (-0.5, -2*\y) --++ (10.75,0)  node[right] {\y$t_{\ell}$};
					}
					% red lines for bounces
					\draw[thick, red,->] (1.5,0) --++ (8.5,-2);% node[above] {$\myv{0}^+$}; 
					\draw[thick, red,->] (10,-2) --++ (-8.5,-2);%node[above] {$\myv{0}^- = \myv{0}^+\Gamma_L$};
					\draw[thick, red,->] (1.5,-4) --++ (8.5,-2);% node[above] {$\myv{0}^{-} \Gamma_s = \myv{0}^+\Gamma_L \Gamma_s$};
				\end{circuitikz}
				} % end scalebox
				\end{center}
			\end{textblock}
\end{frame}

%%%%%%%%%%%%%%%%%%%%%%%%%%%%%%%%%%%%%%%%%%%%%%%%%%%%%
%%%%%%%%%%%%%%%%%%%%%%%%%%%%%%%%%%%%%%%%%%%%%%%%%%%%%

\begin{frame}{Transient Example: Pulsed Source}
			\begin{textblock}{14.5}(0.5,2)
			\begin{center}
				\begin{circuitikz}[scale=0.75]
					\draw (1.5,0) to[open, *-*] ++(0,2)
										to[short, -*] ++(8.5,0);
					\draw (10,0) to[short,*-] ++(-10,0);
					\draw (0,0) to[short] ++(-2,0)
										to[V,a^=10\volts, invert] ++(0,2)
										to[R=25\ohms] ++(2,0)
										to[closing switch, a^=\scriptsize{$0 < t < 2\mu $s}] ++(2,0); % source
					\draw (10,2) to[short] ++(1,0)
										to[R=150\ohms] ++(0,-2)
										to[short] ++(-1,0); % load
					\draw (1.5,3) --++ (8.5,0) node[above, midway] (5,3.25) {$\ell=400$m};
					\draw(1.5,2.75) --++ (0,0.5);
					\draw(10,2.75) --++ (0,0.5);
					\draw[fill=gray!50!white, thick] (2,1.75) rectangle ++(7,0.5); % top tx line
					\draw[fill=gray!50!white, thick] (2,-0.25) rectangle ++(7,0.5); % bottom tx line
					\draw[->] (1.5,2.25) --++(0,-9) ; % left boundary
					\draw[->] (10,2.25) --++(0,-9); % right boundary
					\node[align=left] at (5,1) {\myz{0}=50\ohms \\ $u_p = 0.67c$};
					
					\foreach \y in {1,2,3}{
						\draw[dashed] (-0.5, -2*\y) --++ (10.75,0)  node[right] {\y$t_{\ell}$};
					}
					% red lines for bounces
					\draw[thick, red,->] (1.5,0) --++ (8.5,-2);% node[above] {$\myv{0}^+$}; 
					\draw[thick, red,->] (10,-2) --++ (-8.5,-2);%node[above] {$\myv{0}^- = \myv{0}^+\Gamma_L$};
					\draw[thick, red,->] (1.5,-4) --++ (8.5,-2);% node[above] {$\myv{0}^{-} \Gamma_s = \myv{0}^+\Gamma_L \Gamma_s$};
				\end{circuitikz}
				\end{center}
			\end{textblock}
\end{frame}

%%%%%%%%%%%%%%%%%%%%%%%%%%%%%%%%%%%%%%%%%%%%%%%%%%%%%
%%%%%%%%%%%%%%%%%%%%%%%%%%%%%%%%%%%%%%%%%%%%%%%%%%%%%

\begin{frame}{Transient Example: Pulsed Source}
			\begin{textblock}{8}(0.5,2)
			\begin{center}
			\scalebox{0.5}{
				\begin{circuitikz}[scale=0.75]
					\draw (1.5,0) to[open, *-*] ++(0,2)
										to[short, -*] ++(8.5,0);
					\draw (10,0) to[short,*-] ++(-10,0);
					\draw (0,0) to[short] ++(-2,0)
										to[V,a^=10\volts, invert] ++(0,2)
										to[R=25\ohms] ++(2,0)
										to[closing switch, a^=\scriptsize{$0 < t < 2\mu $s}] ++(2,0); % source
					\draw (10,2) to[short] ++(1,0)
										to[R=150\ohms] ++(0,-2)
										to[short] ++(-1,0); % load
					\draw (1.5,3) --++ (8.5,0) node[above, midway] (5,3.25) {$\ell=400$m};
					\draw(1.5,2.75) --++ (0,0.5);
					\draw(10,2.75) --++ (0,0.5);
					\draw[fill=gray!50!white, thick] (2,1.75) rectangle ++(7,0.5); % top tx line
					\draw[fill=gray!50!white, thick] (2,-0.25) rectangle ++(7,0.5); % bottom tx line
					\draw[->] (1.5,2.25) --++(0,-9) ; % left boundary
					\draw[->] (10,2.25) --++(0,-9); % right boundary
					\node[align=left] at (5,1) {\myz{0}=50\ohms \\ $u_p = 0.67c$};
					
					\foreach \y in {1,2,3}{
						\draw[dashed] (-0.5, -2*\y) --++ (10.75,0)  node[right] {\y$t_{\ell}$};
					}
					% red lines for bounces
					\draw[thick, red,->] (1.5,0) --++ (8.5,-2);% node[above] {$\myv{0}^+$}; 
					\draw[thick, red,->] (10,-2) --++ (-8.5,-2);%node[above] {$\myv{0}^- = \myv{0}^+\Gamma_L$};
					\draw[thick, red,->] (1.5,-4) --++ (8.5,-2);% node[above] {$\myv{0}^{-} \Gamma_s = \myv{0}^+\Gamma_L \Gamma_s$};
				\end{circuitikz}
				} % end scalebox
				\end{center}
			\end{textblock}
\end{frame}



